% method.tex
\section{Methodology}
\label{sec:method}

This section presents the KnowFeat framework in detail. We first provide an overview of the pipeline (Section~\ref{subsec:overview}), then describe each component: multi-type knowledge injection (Section~\ref{subsec:knowledge}), the LLM agent with specialized tools (Section~\ref{subsec:agent}), the three-stage verification pipeline (Section~\ref{subsec:verification}), the intelligent search strategy (Section~\ref{subsec:search}), and explainable provenance output (Section~\ref{subsec:provenance_method}).

\subsection{Framework Overview}
\label{subsec:overview}

KnowFeat is a multi-round, knowledge-guided feature engineering framework built on an LLM agent. Given a tabular dataset with baseline features and target labels, along with a domain knowledge base, KnowFeat generates a set of new features such that appending them to the original data improves predictive performance on a target metric (e.g., AUC or Recall).

The overall pipeline operates as follows (see Figure~\ref{fig:framework} for an illustration):

\begin{enumerate}
    \item \textbf{Initialization.} Load the dataset, compute baseline model performance, and calibrate statistical thresholds via noise injection (Section~\ref{subsec:verification}).
    \item \textbf{Knowledge Assembly.} Organize domain knowledge $\mathcal{K}$ into a five-type taxonomy and inject the selected types into the agent's context (Section~\ref{subsec:knowledge}).
    \item \textbf{Multi-Round Generation.} For each round $r = 1, \ldots, R_{\max}$:
    \begin{enumerate}
        \item Allocate per-dimension budgets based on SHAP feedback from the previous round.
        \item For each semantic dimension (Section~\ref{subsec:search}), the LLM agent explores the data, proposes candidate features grounded in knowledge assets, and submits them for verification.
        \item Each candidate passes through the L1$\to$L2$\to$L3 pipeline. Accepted features are added to the baseline; rejected features enter the failure memory.
        \item Forward selection discovers feature combinations that are individually weak but collectively effective.
    \end{enumerate}
    \item \textbf{Convergence.} The process terminates when no new features are accepted for $C$ consecutive rounds (default $C=2$) or the maximum round count $R_{\max}$ is reached.
    \item \textbf{Output.} The accepted feature set $\mathcal{F}^*$, along with executable code and explainable provenance cards for each feature (Section~\ref{subsec:provenance_method}).
\end{enumerate}

\subsection{Multi-Type Knowledge Injection}
\label{subsec:knowledge}

KnowFeat systematically integrates external domain knowledge into the feature generation process. We organize knowledge into five types of increasing specificity, from general metadata to concrete case evidence:

\textbf{(1) Schema Metadata.} The agent receives only the dataset schema: column names, data types, and basic statistics (mean, cardinality, missing rate). No domain knowledge is provided. This serves as the ablation baseline, testing what the LLM can generate from metadata alone.

\textbf{(2) Risk Indicators.} Structured risk indicators derived from regulatory guidelines and industry standards. Each indicator specifies a risk dimension, risk class, relevant data field, and directionality. For the AML domain, we compile 38 indicators (IND\_01 through IND\_38) covering transaction amount thresholds, frequency patterns, temporal anomalies, and network centrality measures. For example, IND\_05 flags accounts whose single-day transaction count exceeds a regulatory threshold, while IND\_15 identifies abnormal transaction modes such as fan-in and many-to-many patterns.

\textbf{(3) Detection Rules.} Operational detection rules that encode specific logic patterns. Each rule includes a rule ID, natural language description, logical expression, confidence level, and a Python implementation hint. We compile 15 rules (RULE\_001 through RULE\_015) from compliance manuals and internal documentation. For example, RULE\_006 encodes the ``scatter-in, gather-out'' layering pattern.

\textbf{(4) Expert Opinions.} Insights from domain expert interviews that capture tacit knowledge (heuristics, edge cases, and domain intuitions) difficult to formalize as rules but valuable for feature design. We collect 12 expert opinions (EXP\_01 through EXP\_12). For example, EXP\_05 notes that ``border region transactions combined with nighttime activity are a hallmark of cross-border laundering.''

\textbf{(5) Court Document Evidence.} Risk factors extracted from 309 real court judgments via a key-clue extraction pipeline. This is the most concrete type, grounding feature generation in case evidence (e.g., pyramid-shaped fund network patterns independently validated by legal proceedings).

The five types can be freely combined via a knowledge filter parameter, enabling ablation study of each type's contribution. In the ``full'' configuration (default), all types are included. Each knowledge item is assigned a unique identifier for provenance tracking.

\subsection{LLM Agent with Specialized Tools}
\label{subsec:agent}

KnowFeat employs an LLM agent (DeepSeek-V4-0324-Flash in our implementation, with temperature=0 for near-deterministic outputs) that interleaves reasoning with tool calls following the ReAct paradigm~\cite{yao2023react}. The agent is equipped with three specialized tools:

\textbf{Tool 1: Data Exploration.} Executes Python code on the dataset in a sandboxed environment to inspect distributions, compute correlations, and identify patterns before proposing features. The tool returns printed output or return values to the agent.

\textbf{Tool 2: Feature Registration.} Registers a candidate feature with mandatory provenance fields:
\begin{itemize}
    \item \texttt{feature\_name}: A descriptive identifier.
    \item \texttt{code}: Executable Python code that computes the feature from raw data.
    \item \texttt{explanation}: A natural language description of what the feature captures.
    \item \texttt{reasoning}: A structured chain linking knowledge assets to the feature design, formatted as ``Step $N$: Because [knowledge ID] indicates [insight], we compute [transformation]\ldots''
    \item \texttt{knowledge\_ids}: A list of specific knowledge asset identifiers (e.g., RULE\_001, IND\_05, EXP\_03) that inspired the feature.
    \item \texttt{computation\_steps}: Step-by-step description from raw fields to the final feature value.
    \item \texttt{data\_fields\_used}: The raw columns consumed by the feature.
\end{itemize}

\textbf{Tool 3: Verification Trigger.} Triggers the L1$\to$L2$\to$L3 verification pipeline on a registered feature and returns the result (accepted/rejected with detailed reasons). The agent must call Feature Registration followed immediately by Verification Trigger for each candidate; batch submission is prohibited to ensure tight feedback loops.

The agent's prompt instructs it to:
(1) reason about which knowledge assets are relevant to the current dimension and budget,
(2) explore the data to validate assumptions,
(3) propose a feature grounded in specific knowledge items, and
(4) submit it for verification and react to the feedback.

\subsection{Three-Stage Verification Pipeline}
\label{subsec:verification}

Each candidate feature passes through three progressively stricter verification stages:

\subsubsection{L1: Code Execution Sandbox}
The generated Python code is executed in an isolated sandbox with a 30-second timeout. L1 checks for:
\begin{itemize}
    \item Syntax and runtime errors.
    \item Output shape consistency (one value per sample).
    \item Presence of required identifier columns.
    \item Data type validity (numeric output).
\end{itemize}
If L1 fails, the agent receives the error message and is allowed up to 2 automatic repair attempts before the feature is permanently rejected.

\subsubsection{L2: Statistical Quality Filtering}
The feature's statistical properties are evaluated against automatically calibrated thresholds. To establish these thresholds for discriminative power, we generate $N_{\text{noise}} = 100$ random features (drawn from standard distributions), compute their quality metrics, and set the acceptance threshold at the $P_{\text{noise}}$-th percentile (default $P_{\text{noise}} = 95$). This ensures that accepted features are statistically better than random noise. L2 checks four criteria:
\begin{itemize}
    \item \textbf{Information Value (IV)}~\cite{chen2018aml}: Must exceed the noise-calibrated threshold, ensuring the feature has meaningful discriminative power between positive and negative classes.
    \item \textbf{Kolmogorov--Smirnov (KS) Statistic}~\cite{nargesian2017learning}: Must exceed the noise-calibrated threshold, measuring the maximum distributional separation between classes.
    \item \textbf{Redundancy}: Maximum absolute Pearson correlation with all existing accepted features must be below 0.7, preventing near-duplicate features from inflating the feature set.
    \item \textbf{Missing Rate}: The proportion of NaN values must be at most 50\%, ensuring sufficient coverage for model training.
\end{itemize}
The noise-calibrated thresholds apply to IV and KS; the redundancy and missing rate thresholds are fixed.

\subsubsection{L3: Model Effectiveness}
The feature is evaluated for its actual contribution to predictive performance. We train an XGBoost~\cite{chen2016xgboost} model on the baseline features augmented with the candidate feature, using stratified 80/20 train-test splits across 5 random seeds $\{42, 123, 456, 789, 2024\}$. The acceptance criterion depends on the task:
\begin{itemize}
    \item \textbf{Classification (public datasets)}: $\Delta\text{AUC} > 0$, i.e., the augmented model must achieve strictly higher AUC than the baseline.
    \item \textbf{AML detection}: $\Delta\text{Recall} > 0$ \emph{and} Precision $\geq 90\%$, ensuring that the feature improves detection without sacrificing precision.
\end{itemize}

The three-stage design is intentionally progressive: L1 is cheap (milliseconds), L2 is moderate (seconds), and L3 is expensive (minutes per feature). By filtering early, we avoid wasting computational resources on clearly defective candidates.

\subsection{Intelligent Search Strategy}
\label{subsec:search}

To systematically explore the feature space, KnowFeat employs a multi-component search strategy.

\textbf{Dimension-Based Budgeting.}
We decompose the feature space into $D$ semantic dimensions relevant to the target domain (default $D=4$ for the AML domain): Amount \& Frequency, covering monetary values and transaction counts; Time \& Geography, capturing temporal patterns and spatial distributions; Behavior Sequence, encoding sequential patterns in user actions; and Entity Association, derived from relationships between entities such as accounts and counterparties.

Each dimension receives an initial budget of $B_0$ candidates per round (default $B_0 = 3$). After each round, SHAP~\cite{lundberg2017shap} (SHapley Additive exPlanations) analysis is performed on the current model to estimate each dimension's contribution to predictions.

\textbf{SHAP-Guided Budget Reallocation.}
Let $s_j$ denote the mean absolute SHAP value for dimension $j$'s accepted features, $B_{\min}$ the minimum per-dimension budget (default 1), and $B_{\text{total}}$ the total budget for the round. We reallocate the budget for round $r{+}1$ using an \emph{inverse-SHAP weighting} strategy:
\begin{equation}
    B_j^{(r+1)} = B_{\min} + (B_{\text{total}} - D \cdot B_{\min}) \cdot \frac{1/s_j}{\sum_{j'=1}^{D} 1/s_{j'}}
\end{equation}
This ensures that underrepresented dimensions (those contributing least to the current model) receive more exploration budget, promoting balanced feature coverage.

\textbf{Failure Memory.}
A failure memory $\mathcal{M}_{\text{fail}}$ records all rejected features along with their rejection reasons. Before each generation step, the $K$ most recent failures for the current dimension (default $K=10$) are injected into the agent's prompt, preventing the regeneration of previously rejected ideas and guiding the agent toward novel directions.

\textbf{Forward Selection.}
Individual features may fail L3 yet contribute positively when combined with others. We apply greedy forward selection on L2-passed but L3-failed candidates, sorted by L2 quality scores. Each candidate is added to the accepted set if the augmented set improves the target metric, capturing synergistic feature interactions.

\textbf{Cross-Dimension Propagation.}
After each round, the top-$K_{\text{cross}}$ features (default $K_{\text{cross}} = 5$) by SHAP importance are shared across all dimensions as ``inspiration prompts,'' encouraging the agent to adapt successful patterns from one dimension to others.

\subsection{Explainable Provenance Output}
\label{subsec:provenance_method}

Every generated feature is accompanied by a structured \emph{provenance card} that enables full explainability. Each provenance card records:
\begin{itemize}
    \item \textbf{Knowledge assets}: The specific identifiers (e.g., RULE\_006, IND\_15, EXP\_05) that inspired the feature design.
    \item \textbf{Reasoning chain}: A step-by-step explanation of how domain knowledge led to the feature, formatted as ``Step $N$: Because [knowledge ID] indicates [insight], we compute [transformation]\ldots''
    \item \textbf{Computation steps}: The transformation pipeline from raw data fields to the final feature value.
    \item \textbf{Verification results}: L1/L2/L3 outcomes including IV, KS, redundancy, and $\Delta$AUC/$\Delta$Recall.
\end{itemize}
This traceability is essential for regulated domains: compliance officers can trace \emph{why} each feature exists, \emph{what evidence} supports it, and \emph{how} it was verified. Section~\ref{sec:case_study} illustrates provenance cards through concrete examples.
